\scanpage{70}\mspage{E.}
Le deuxième est est appelé le \uline{groupe fondamental
birationnel} de $K/k$. [N.B. il y aurait
\unsure{à} l'introduire plutôt comme groupe alg.
sur $K$], \inserted{noté $\pi_1(K/k;\Omega)$. Pour tt modèle \uline{propre} \uline{projectif} $X$, $\pi_1(X;\Omega)$ est un quotient de $\pi_1(K/k)$.} On voit donc qu'il est \unsure{isomorphe|toujours}
\unsure{à un|un} quotient de $\pi_1(U)$. Donc, comme
$U$ peut être pris complémentaire d'une partie fermée
de codim $\geq 2$ ds un \unsure{schéma|sch.} intègre propre :

\uline{Th.} Si $K/k$ est une extension de
type fini, avec $k$ alg. clos, alors
\deleted{\illegible{1}} $\pi_1(K/k;\Omega)$ est top.$^{\mathrm t}$ de type fini.

\margin{\uline{Remarque} Si $K$ a un modèle propre régulier $X$, alors pour $X'$ \deleted{propre} normal dominant $X$, on a $\pi_1(X';\Omega)\overset{\sim}{\longrightarrow}\pi_1(X;\Omega)$, donc le système projectif des $\pi_1(X';\Omega)$, $X'$ propre normal, est essentiellement constant. Qu'en est-il en général ??? Si $X$ est un modèle propre normal, \deleted{\unsure{tout}} $U$ l'ens de ses pts réguliers, est-il vrai que $\pi_1(U;\Omega)\to\pi_1(K/k;\Omega)$ est bijectif, i.e. que tte \deleted{\unsure{variété $X$}} \unsure{extension} finie $K'$ de $K$, non ramifiée sur $U$, est non ramifiée sur un modèle \inserted{\unsure{normal}} $X'$ \unsure{dérivant de|distinct de|dominant} $X$ pour \unsure{certains|certain} \unsure{$X'$|$X$} ? C'est une question en fait de schémas généraux ....}
% slanted
% archive sheet number: 69
